\hypertarget{prescriptive-remediation-and-cicd-quality-gating}{%
\section{Prescriptive Remediation and CI/CD Quality Gating}\label{prescriptive-remediation-and-cicd-quality-gating}}

\label{sec:6}

Attribution (Section~\ref{sec:4}) and impact analysis (Section~\ref{sec:5}) are diagnostic: they tell
an architect \emph{which} components to harden and \emph{why}. This section closes the loop with a
prescriptive stage that proposes concrete architectural edits and verifies that they reduce simulated
failure impact before any deployment, and describes how the diagnostics are operationalised as a
continuous CI/CD quality gate.

\hypertarget{generateverify-and-the-remediation-operators}{%
\subsection{Generate--Verify and the Remediation Operators}\label{generateverify-and-the-remediation-operators}}

\label{sec:6.1}

Remediation runs in two strictly separated phases, preserving the independence discipline of the rest
of the framework. \textbf{Generate}: given \(G_{\text{analysis}}\) and its attribution, four
operators propose candidate topology edits from structure alone, never reading simulated impact:

\begin{longtable}[]{@{}
  >{\raggedright\arraybackslash}p{(\columnwidth - 6\tabcolsep) * \real{0.25}}
  >{\raggedright\arraybackslash}p{(\columnwidth - 6\tabcolsep) * \real{0.25}}
  >{\raggedright\arraybackslash}p{(\columnwidth - 6\tabcolsep) * \real{0.28}}
  >{\raggedright\arraybackslash}p{(\columnwidth - 6\tabcolsep) * \real{0.22}}@{}}
\caption{The four remediation operators, their structural triggers, and the failure mode each targets.\label{tab:4}}\tabularnewline
\toprule
Operator & Structural trigger & Edit applied & Failure mode targeted \\
\midrule
\endfirsthead
\toprule
Operator & Structural trigger & Edit applied & Failure mode targeted \\
\midrule
\endhead
\textbf{RedundancyInsertion} & directed articulation point / high \(A\) SPOF & add a redundant instance or redistribute responsibilities & graph-partitioning SPOF \\
\textbf{PathDiversification} & bridge edge / single routing path for a topic & add an alternative route (e.g.~a second routing broker or network link) & fragmentation on a non-redundant edge \\
\textbf{FanOutReduction} & high structural blast radius (topic subscriber fan-out; library consumer count) & interpose an intermediary or split the over-shared channel & simultaneous blast / fan-out explosion \\
\textbf{SharedTopicReduction} & high multi-path coupling (large \texttt{path\_count} / MPCI between a pair) & decouple redundant shared topics between the pair & multi-channel coupling fragility \\
\bottomrule
\end{longtable}

\textbf{Verify}: each candidate edit is applied to produce a counterfactual graph, on which the
\texttt{FailureSimulator} (Section~\ref{sec:5.1}) is re-run from scratch; the edit is accepted only if

\begin{equation}
\Delta I > \kappa\,\sigma_{\text{seed}}
\end{equation}

--- the reduction in \(I_{\text{comp}}\) exceeds \(\kappa = 1.0\) times the simulator's own
across-seed noise --- at \emph{every} sampled propagation threshold, evaluated per candidate edit
before it is committed to the policy. This separation matters: a stage that both proposed and scored
edits using the same signal would be optimizing against itself, and normalising the acceptance bar by
\(\sigma_{\text{seed}}\) ties it to the fragility of the cascade at that point rather than to an
arbitrary constant. FanOutReduction's trigger is deliberately \emph{not} the composite \(Q(v)\) but a
direct structural blast-radius signal (subscriber fan-out, consumer count), so that a low-\(Q\),
high-blast component is still selected --- the remediation-side expression of the paper's central
claim that single-score criticality is insufficient, independent of whether Section~\ref{sec:5.4}
finds such a component in this particular suite. No Verify result feeds back into Generate within a
run, so there is no closed-loop search that would reintroduce the circularity the framework is built
to avoid.

Running this procedure across the scenario suite (\(\kappa = 1.0\), three propagation thresholds,
three seeds --- reduced from five because acceptance requires a full re-simulation per candidate per
threshold, the most expensive experiment in the study), 162 of 332 candidate edits (48.8\%) clear the
filter, with no scenario in the suite regressing under the per-edit filter, removing by construction
the failure mode of an earlier, unverified design in which a regressing edit could be carried by an
improving one. Acceptance is decided on singletons, not subsets: nothing in the procedure establishes
that a set of individually-accepted edits remains beneficial applied together. The acceptance rate
varies widely with topology --- 3 of 35 candidates on Autonomous Vehicle against 38 of 58 on IoT
Smart City --- concentrated where a fan-out bottleneck actually exists, and the resulting System Risk
Index reductions are real but small across the suite (\(+0.0025\) to \(+0.0158\)), which is the
result we report rather than a demonstration that the prescriptive stage is yet practically valuable
(Section~\ref{sec:9.3}).

\hypertarget{the-cicd-quality-gate}{%
\subsection{The CI/CD Quality Gate}\label{the-cicd-quality-gate}}

\label{sec:6.2}

To operationalise these diagnostics, SaG integrates into developer workflows as a blocking Quality
Gate: when a pull request introduces Architecture-as-Code changes, the pipeline executes
\texttt{detect\_antipatterns.py}, which runs the full anti-pattern catalog against the candidate
topology and issues an exit code --- 0 (no anomalies, deployment permitted), 1 (medium-severity
smells, permitted with warnings), or 2 (CRITICAL/HIGH severity anomalies --- SPOFs, cyclic
dependencies, broker overload --- build broken, \textbf{deployment blocked}). The underlying
machinery is in-memory via the thread-safe \texttt{MemoryRepository} port and does not require a live
database connection, which is what makes the runtimes of Section~\ref{sec:8.4} possible.

This is an \emph{absolute} gate: every run evaluates the candidate topology's full finding set, not a
diff against a prior baseline. This has a known consequence we do not paper over: a real architecture
carrying an intentional, risk-accepted single point of failure --- a sole-source feed, a deliberately
unreplicated legacy broker --- fails the build on every commit, indistinguishable from a genuine
regression. A \emph{delta-aware} gate that evaluates the candidate against a merge-base topology,
blocking only on newly introduced findings, together with a waiver register recording accepted risk
so it stays visible rather than silently re-triggering, would close this gap; we describe the design
in Section~\ref{sec:9.3} as future work rather than claim it here, since the mechanism is not
implemented in the released tool.

\begin{center}\rule{0.5\linewidth}{0.5pt}\end{center}
